\chapter{Research Workflows (End-to-end Loops)}

\ChapterMeta{
This chapter defines repeatable research loops and the utilities for subsets, sweeps, and benchmarks.
}{
It keeps experiments artifact-first so results can be compared and regressed over time.
}{
It assumes stable JSON artifacts and, when needed, protocol presets. Keep subsets deterministic so comparisons stay fair and valid.
}
\ChapterAudience{Read this chapter when you are planning ablations, sweeps, offline distillation, or artifact-first research loops.}

\WorkflowOpening
{Run repeatable research loops while keeping the evaluation boundary artifact-first and comparable.}
{Use this for ablations, frozen subsets, protocol-pinned comparisons, sweeps, prediction distillation, and research notes.}
{Stable JSON artifacts, a fixed dataset subset or protocol, and explicit output paths for reports and histories.}
{Start with validation and protocol-pinned eval; use the specific sweep, distillation, or benchmark helper only after artifacts are valid.}
{Suite reports, sweep histories, distilled predictions, benchmark summaries, and research-note inputs.}
{Check that subset seeds, protocol files, metric keys, and artifact paths are recorded before interpreting the result.}
{Changing subsets mid-study, mixing conceptual and executable commands, skipping validation, or comparing reports with different protocol settings.}

\section{The research loop in \yolozu{}}
A canonical research iteration proceeds as follows:
\begin{enumerate}
  \item Build or acquire a model (either via internal training or external sourcing).
  \item Export predictions that follow the stable JSON interface contract.
  \item Validate artifacts immediately and fail fast.
  \item Evaluate the artifacts under a pinned protocol (generating suite reports).
  \item Run backend parity checks and compare historical runs with JSONL/CSV/MD histories.
  \item (Optional) Apply post-hoc calibration and then re-evaluate.
  \item (Optional) Execute hyperparameter sweeps and benchmark latency metrics.
\end{enumerate}

The discipline is simple: \textbf{treat artifacts as first-class entities}. Every run must generate JSON files that support reproduction and comparison.

\section{Freezing a dataset subset for ablations}
Rapid ablation studies frequently require a deterministic subset of a larger dataset. \yolozu{} provides a dedicated utility for this purpose:
\begin{lstlisting}[language=bash]
python3 tools/make_subset_dataset.py   --dataset data/coco128   --split train2017   --n 50   --seed 0   --out reports/coco128_50
\end{lstlisting}

This command writes:
\begin{itemize}
  \item A dedicated subset dataset directory.
  \item A frozen image list for reproducible sampling.
  \item A subset metadata JSON file for clear provenance.
\end{itemize}

\section{Protocol-pinned evaluation (YOLO26 example)}
Protocols pin evaluation settings so metrics remain comparable over time. For YOLO26, the protocol is defined in \path{protocols/yolo26_eval.json}; the evaluation tool reads this file and embeds its parameters into the final suite reports.

To execute a suite evaluation across multiple prediction artifacts:
\begin{lstlisting}[language=bash]
python3 tools/eval_suite.py   --protocol yolo26   --dataset /path/to/yolo-dataset   --predictions-glob '/path/to/pred_*.json'   --output reports/eval_suite.json
\end{lstlisting}

The resulting report includes protocol metadata, such as split definitions and bounding box formats. This metadata is required for valid research comparisons.

\section{Validating interface contracts as a research quality gate}
Before investing time in metric interpretation, you must validate your artifacts:
\begin{lstlisting}[language=bash]
python3 tools/validate_dataset.py --dataset /path/to/yolo --strict
python3 tools/validate_predictions.py /path/to/predictions.json --strict
\end{lstlisting}

This step is particularly critical when ingesting predictions exported from external pipelines.

\section{Backend parity (Torch vs. ONNX Runtime vs. TensorRT)}
When evaluating different inference backends, output alignment comes first. Use the parity tooling to compare two \path{predictions.json} artifacts and identify systematic differences, such as preprocessing, coordinate formatting, or numerical drift.

\textbf{Conceptual (not copy-paste as-is):} parity command shape.
Use real artifacts and validate them first as shown in the quality-gate section above.
\begin{lstlisting}[language=bash]
yolozu parity   --reference reports/pred_torch.json   --candidate reports/pred_onnxrt.json
\end{lstlisting}

\section{Sweeps (Hyperparameter search harness)}
Research workflows often need parameter sweeps that run outside the primary training loop. The sweep harness runs external commands, aggregates metrics, and writes summary tables.

Quick start:
\begin{lstlisting}[language=bash]
python3 tools/hpo_sweep.py --config docs/hpo_sweep_example.json --resume
\end{lstlisting}

Default outputs include:
\begin{itemize}
  \item \path{reports/hpo_sweep.jsonl}
  \item \path{reports/hpo_sweep.csv}
  \item \path{reports/hpo_sweep.md}
\end{itemize}

Use the \cmd{--resume} flag to make sweeps incremental and resilient to interruptions.

Critical configuration fields (detailed in \path{docs/hpo_sweep.md}):
\begin{itemize}
  \item \textbf{Required}: \cmd{base_cmd}, alongside either \cmd{param_grid} or \cmd{param_list}.
  \item \textbf{Optional}: \cmd{run_dir}, \cmd{metrics.path}, \cmd{metrics.keys}.
  \item \textbf{Optional outputs}: \cmd{result_jsonl}, \cmd{result_csv}, \cmd{result_md}.
  \item \textbf{Optional runtime controls}: \cmd{env}, \cmd{shell}.
\end{itemize}

The harness is intentionally framework-agnostic; it makes no assumptions about the underlying training stack.

\section{Prediction distillation workflow}
\yolozu{} includes a beginner-friendly \emph{offline prediction distillation} helper at
\path{tools/distill_predictions.py}. This section is intentionally written in the same spirit as the TTT chapter:
start from the goal, then follow the concrete procedure, then inspect the resulting artifacts.

\subsection{What this workflow is}
Prediction distillation in this chapter means:
\begin{itemize}
  \item take a \textbf{student} \path{predictions.json},
  \item take a \textbf{teacher} \path{predictions.json},
  \item compare them image-by-image,
  \item blend matched detections, optionally inject selected teacher-only detections,
  \item write a new distilled \path{predictions.json} plus a compact report.
\end{itemize}

The important point is that this is an \textbf{artifact-level} workflow.
It does \emph{not} update model weights. It is a fast research helper for answering:
\begin{quote}
``If I trust this teacher more than the student on this subset, is there enough signal to justify a heavier training-time distillation setup?''
\end{quote}

\subsection{What this workflow is not}
This section's \emph{prediction distillation} is narrower than the self-distillation discussed in the continual-learning chapter and narrower than TTT:
\begin{itemize}
  \item it is \textbf{not} checkpoint-based self-distillation during training,
  \item it is \textbf{not} an anti-forgetting method by itself,
  \item it is \textbf{not} changing model parameters at inference time,
  \item it is \textbf{not} a substitute for the full evaluator when reporting final numbers.
\end{itemize}

So the separation is:
\begin{itemize}
  \item \textbf{Prediction distillation}: rewrite one prediction artifact using another prediction artifact.
  \item \textbf{Self-distillation}: train a checkpoint while regularizing it toward another checkpoint.
  \item \textbf{TTT / CTTA}: adapt model parameters in memory during inference-time operation.
\end{itemize}

\begin{figure}[tbp]
\centering
\resizebox{0.94\linewidth}{!}{%
\begin{tikzpicture}[
  font=\small,
  box/.style={draw, rounded corners, align=left, inner sep=6pt},
  artifactbox/.style={box, fill=blue!8, text width=3.8cm},
  matchbox/.style={box, fill=orange!10, text width=4.4cm},
  actionbox/.style={box, fill=green!10, text width=3.8cm},
  outputbox/.style={box, fill=purple!8, text width=4.0cm},
  arrow/.style={-{Latex[length=2mm]}, thick}
]

\node[artifactbox] (student) at (0,2.45) {\textbf{Student artifact}\\\path{predictions_student.json}};
\node[artifactbox] (teacher) at (9.2,2.45) {\textbf{Teacher artifact}\\\path{predictions_teacher.json}};
\node[matchbox] (match) at (4.6,0.15) {\textbf{Match detections}\\same image, same class, IoU threshold};
\node[actionbox] (blend) at (0,-2.75) {\textbf{Blend matched scores}\\controlled by \cmd{--alpha}};
\node[actionbox] (inject) at (9.2,-2.75) {\textbf{Optionally inject teacher-only detections}\\guarded by score / duplicate / count caps};
\node[outputbox] (output) at (4.6,-5.55) {\textbf{Distilled artifacts}\\\path{predictions_distilled.json}\\\path{distill_report.json}};

\draw[arrow] (student) -- (match);
\draw[arrow] (teacher) -- (match);
\draw[arrow] (match) -- (blend);
\draw[arrow] (match) -- (inject);
\draw[arrow] (blend) -- (output);
\draw[arrow] (inject) -- (output);

\end{tikzpicture}
}
\caption{Prediction distillation flow in \yolozu{}: keep the student artifact as the base, borrow teacher signal under explicit rules, then emit a new distilled artifact instead of mutating the originals.}
\end{figure}

\subsection{A friendlier mental model}
If you need to explain this workflow to a new teammate, the most useful sentence is:
\begin{quote}
The student artifact is the first draft; the teacher artifact is the reviewer; distillation applies only a limited set of corrections and writes a new artifact.
\end{quote}

In practice, each candidate object lands in one of three buckets:
\begin{itemize}
  \item \textbf{Matched}: both student and teacher found the same object. This usually leads to a safer score update.
  \item \textbf{Teacher-only}: the teacher found an object that the student missed. This can improve recall, but it is also where false positives enter.
  \item \textbf{Student-only}: the student found an object that the teacher did not. The helper keeps the student artifact as the base.
\end{itemize}

\begin{figure}[h]
\centering
\resizebox{\linewidth}{!}{%
\begin{tikzpicture}[
  font=\small,
  box/.style={draw, rounded corners, align=left, inner sep=6pt, text width=0.28\linewidth},
  arrow/.style={-{Latex[length=2mm]}, thick}
]

\node[box] (matched) at (0,0) {\textbf{Matched case}\\Student and teacher overlap on the same object.\\Most common effect: confidence blending.\\Typical question: ``Is the teacher just better calibrated?''};
\node[box] (teacher) at (6.3,0) {\textbf{Teacher-only case}\\Teacher sees an object the student missed.\\Most common effect: recall gain or extra noise.\\Typical question: ``Is this a real missing object?''};
\node[box] (student) at (12.6,0) {\textbf{Student-only case}\\Student predicts something the teacher does not.\\Most common effect: artifact stays unchanged.\\Typical question: ``Is the teacher blind here or is the student wrong?''};

\draw[arrow] (matched.south) -- ++(0,-0.9) node[below, align=center] {\textbf{Safer path}\\small score correction};
\draw[arrow] (teacher.south) -- ++(0,-0.9) node[below, align=center] {\textbf{Highest-variance path}\\possible recall gain};
\draw[arrow] (student.south) -- ++(0,-0.9) node[below, align=center] {\textbf{Base preserved}\\no teacher intervention};

\end{tikzpicture}
}
\caption{The three cases to look for when reading a distillation result. For beginners, this is often easier to reason about than starting directly from the final metric table.}
\end{figure}

\subsection{Step-by-step procedure}
\paragraph{Step 1: Prepare two valid artifacts}
You need one student predictions artifact and one teacher predictions artifact.
Before distilling, validate both:
\begin{lstlisting}[language=bash]
python3 tools/validate_predictions.py reports/predictions_student.json --strict
python3 tools/validate_predictions.py reports/predictions_teacher.json --strict
\end{lstlisting}

For operator convenience, \cmd{tools/distill_predictions.py} also accepts wrapped payloads with partial \path{meta} blocks and supplements missing stable keys (for example \texttt{adapter}, \texttt{timestamp}, \texttt{tta}, \texttt{ttt}) before validation. This makes checked-in research artifacts easier to reuse without hand-editing metadata first.

\paragraph{Step 2: Decide how aggressive the blend should be}
The most important knobs are:
\begin{itemize}
  \item \cmd{--alpha}: how strongly matched detections move toward the teacher score.
  \item \cmd{--iou-threshold}: how much overlap teacher/student detections need before they count as the same object.
  \item \cmd{--add-missing}: whether teacher-only detections may be injected.
  \item \cmd{--teacher-min-score}: minimum confidence required before a teacher-only detection may be injected.
  \item \cmd{--max-added-per-image}: cap on teacher-only additions.
  \item \cmd{--add-duplicate-iou-threshold}: prevents near-duplicate teacher injections.
\end{itemize}

\paragraph{Safe-first reading strategy}
When in doubt, use this order:
\begin{enumerate}
  \item run a conservative pass first,
  \item inspect whether \path{matched} dominates \path{added},
  \item only then enable teacher-only injection,
  \item if \path{added} dominates but the visual result looks noisy, treat the gain as fragile.
\end{enumerate}

\paragraph{Step 3: Run a conservative first pass}
Start with a conservative run before enabling teacher-only injection:
\begin{lstlisting}[language=bash]
python3 tools/distill_predictions.py \
  --student reports/predictions_student.json \
  --teacher reports/predictions_teacher.json \
  --dataset data/coco128 \
  --split val2017 \
  --alpha 0.5 \
  --iou-threshold 0.7 \
  --output reports/predictions_distilled.json \
  --output-report reports/distill_report.json
\end{lstlisting}

\paragraph{Step 4: Run an exploratory ablation if needed}
If you want to test whether teacher-only detections help, enable guarded injection:
\begin{lstlisting}[language=bash]
python3 tools/distill_predictions.py \
  --student reports/predictions_student.json \
  --teacher reports/predictions_teacher.json \
  --dataset data/coco128 \
  --split val2017 \
  --alpha 0.5 \
  --iou-threshold 0.7 \
  --add-missing \
  --teacher-min-score 0.25 \
  --max-added-per-image 20 \
  --add-duplicate-iou-threshold 0.9 \
  --output reports/predictions_distilled.json \
  --output-report reports/distill_report.json
\end{lstlisting}

\paragraph{Step 4b: Prefer a checked-in YAML config when the CLI gets too long}
The helper also accepts \cmd{--config} in JSON or YAML format.
For day-to-day use, the shortest recommended boilerplate is the checked-in YAML file:
\path{configs/examples/distill_predictions.yaml}.

\begin{lstlisting}[language=bash]
python3 tools/distill_predictions.py \
  --student reports/predictions_student.json \
  --teacher reports/predictions_teacher.json \
  --dataset data/coco128 \
  --split val2017 \
  --config configs/examples/distill_predictions.yaml \
  --output reports/predictions_distilled.json \
  --output-report reports/distill_report.json
\end{lstlisting}

This is the preferred operator path when the long explicit flag list becomes hard to remember or copy correctly.

\paragraph{Step 5: Read the outputs in the right order}
The recommended reading order is:
\begin{enumerate}
  \item \path{distill_report.json}
  \item \path{predictions_distilled.json}
  \item a full evaluator report if the result looks promising
\end{enumerate}

\subsection{How the algorithm behaves}
\paragraph{Matched detections}
When teacher and student detections overlap enough under \cmd{--iou-threshold}, the helper treats them as corresponding detections and blends scores using \cmd{--alpha}.
This is usually the safer path because both artifacts already agree that there is probably an object there.

\paragraph{Teacher-only detections}
With \cmd{--add-missing}, the helper may inject detections that only the teacher predicted.
This is where gains and risks both tend to come from.

Typical upside:
\begin{itemize}
  \item recover missed objects,
  \item reveal useful teacher signal quickly.
\end{itemize}

Typical downside:
\begin{itemize}
  \item propagate teacher false positives,
  \item grow duplicate boxes when guardrails are too loose,
  \item overstate benefit on a tiny subset.
\end{itemize}

\subsection{How to read the report}
The compact report is enough for first-pass interpretation.
Inspect:
\begin{itemize}
  \item \path{losses.distill_score_gap}
  \item \path{metrics.student.map50}
  \item \path{metrics.student.map50_95}
  \item \path{metrics.distilled.map50}
  \item \path{metrics.distilled.map50_95}
  \item \path{meta.matched}
  \item \path{meta.added}
  \item \path{meta.distill.*}
\end{itemize}

Interpretation hints:
\begin{itemize}
  \item \textbf{high matched, low added}: mostly confidence refinement on objects both models found.
  \item \textbf{high added}: the result depends heavily on teacher-only injection.
  \item \textbf{metrics up with low added}: usually a healthier sign.
  \item \textbf{metrics up with very high added}: inspect duplicates and false positives carefully.
\end{itemize}

\subsection{Pros and Cons}
\begin{longtable}{@{}p{0.22\textwidth}p{0.33\textwidth}p{0.33\textwidth}@{}}
\toprule
\textbf{Workflow} & \textbf{Pros} & \textbf{Cons} \\
\midrule
Prediction distillation & Fast CPU-friendly ablation; no retraining required; preserves the predictions interface contract workflow; easy side-by-side artifact comparison & Post-hoc artifact rewrite only; gains can be fragile when injection dominates; does not solve forgetting by itself \\
\bottomrule
\end{longtable}

\subsection{Practical guardrails}
Recommended defaults for safer use:
\begin{itemize}
  \item keep \cmd{--teacher-min-score >= 0.2},
  \item keep \cmd{--max-added-per-image} finite,
  \item keep \cmd{--add-duplicate-iou-threshold} high (for example, \texttt{0.85}--\texttt{0.95}),
  \item specify \cmd{--split} when the dataset root has multiple candidate splits,
  \item treat the lightweight \cmd{simple_map} result as a \emph{quick signal}, not the final claim.
\end{itemize}

\subsection{When to use prediction distillation vs TTT vs self-distillation}
\begin{longtable}{@{}p{0.30\textwidth}p{0.56\textwidth}@{}}
\toprule
\textbf{Goal} & \textbf{Recommended workflow} \\
\midrule
Quick offline teacher/student ablation & Prediction distillation \\
Adapt at inference time under domain shift & TTT / CTTA \\
Mitigate forgetting while training across tasks/domains & Continual learning with self-distillation / replay / PEFT \\
Learn a new checkpoint that actually absorbs teacher behavior & Training-time distillation \\
\bottomrule
\end{longtable}

\subsection{Background references}
This helper is an artifact-level utility, not a claim of faithful reproduction of any single paper.
The conceptual lineage still comes from the broader distillation literature:
\begin{itemize}
  \item knowledge distillation: \cite{hinton2015distilling},
  \item self-distillation / Born-Again Networks: \cite{furlanello2018bornagain}.
\end{itemize}

Within research loops, prediction distillation is most useful as a \textbf{fast feasibility check}:
keep the dataset subset fixed, keep the downstream evaluator fixed, and use the distilled artifact to answer whether there is teacher signal worth pursuing further.

\section{Latency/FPS benchmarks with historical tracking (JSONL)}
Benchmark reports must be treated as formal artifacts. The latency harness generates a stable JSON report and can append results to a JSONL history file:
\begin{lstlisting}[language=bash]
python3 tools/benchmark_latency.py   --iterations 200   --warmup 20   --output reports/benchmark_latency.json   --history reports/benchmark_latency.jsonl   --notes 'baseline on target HW'
\end{lstlisting}

This dual-output pattern (a single JSON report coupled with a JSONL history) is invaluable for systematically tracking performance regressions across successive commits and varying hardware configurations.

\section{Run contract artifacts (Training research)}
When utilizing the in-repository trainer, it is strongly recommended to employ the contracted run mode to ensure the generation of consistent artifacts. This methodology is comprehensively detailed in Chapter~7 (Training and the Run Contract).

Minimal example:
\begin{lstlisting}[language=bash]
yolozu train configs/examples/train_contract.yaml --run-id exp01
\end{lstlisting}

Contracted outputs are systematically organized under \path{runs/exp01/} and encompass checkpoints, metrics JSONL files, ONNX exports, and parity reports.

\section{Notes on long-tail calibration experiments}
Post-hoc calibration frequently serves as a primary research axis. Utilize \cmd{yolozu calibrate} to generate calibrated prediction artifacts, enabling rigorous side-by-side comparisons of different methodologies. This process is thoroughly summarized in Chapter~6 (Score Calibration and Long-tail Adjustments).
